% TOP_PROBLEM: {"id": 20000897, "title": "Popular-direction planarization for no-five-coplanar sumsets", "category_id": 2, "category": {"id": 2, "name": "combinatorics", "display_name": "Combinatorics", "description": "Counting problems, graph theory, discrete structures.", "slug": "combinatorics", "order_index": 2, "created_at": "2026-07-31T15:26:25.671Z"}, "status": "partially_solved", "problem_number": "AIM-COMBINATORICS-0022", "set_id": 14, "statusQualification": "Uses the precise three-dimensional extremal target; omits the source background exponent typo in its comparison with planar examples."}
% FIRST_POSTED: 2026-10-01
% ENTRY_KIND: statement_only
% UnsolvedMath ID 20000897; source catalogue code AIM-COMBINATORICS-0022
% !TeX program = lualatex
% Definition and sources only; this file is not an LLM solution attempt.
\documentclass[11pt,a4paper]{article}
\usepackage[margin=25mm]{geometry}
\usepackage{fontspec}
\setmainfont{lmroman10-regular.otf}[BoldFont=lmroman10-bold.otf,
 ItalicFont=lmroman10-italic.otf,BoldItalicFont=lmroman10-bolditalic.otf]
\usepackage{amsmath,amssymb,mathtools}
\usepackage{unicode-math}
\setmathfont{latinmodern-math.otf}
\AtBeginDocument{\let\setminus\smallsetminus}
\usepackage{xurl,hyperref}
\hypersetup{colorlinks=true,urlcolor=blue}
\setcounter{secnumdepth}{0}
\setlength{\parindent}{0pt}
\setlength{\parskip}{5pt}
\setlength{\emergencystretch}{3em}
\begin{document}
\section*{Popular-direction planarization for no-five-coplanar sumsets}\label{20000897}
{\small UnsolvedMath ID 20000897 · AIM-COMBINATORICS-0022 · Combinatorics · AIM Workshop Problem Lists}
\subsection{Definitions and mathematical statement}
For a finite subset $A\subset\mathbb R^3$, define its sumset by
$A+A=\{a+b:a,b\in A\}$, using ordinary vector addition and allowing $a=b$.
Different ordered pairs yielding the same vector contribute only one sumset element.
Say that $A$ satisfies the plane condition if every affine plane in $\mathbb R^3$ contains at most four distinct points of $A$.
The condition concerns all affine planes, including planes not passing through the origin.

For each integer $n\geq5$, set
\[
 D_3(n)=\min\{|A+A|: A\subset\mathbb R^3,
       |A|=n,\ A\text{ satisfies the plane condition}\}.
\]
Such configurations exist, and the minimum is meaningful because the possible sumset sizes are integers in a finite range.
The problem is to determine the asymptotic order of $D_3(n)$, or equivalently the largest universal growth factor $f(n)=D_3(n)/n$ in the inequality $|A+A|\geq n f(n)$ for these configurations.
A sharp result should supply both a bound valid for every admissible set and constructions showing how small the doubling ratio can be.
In particular, the requested factor depends only on cardinality; dependence on the individual set would not give a uniform extremal statement.
No separation, bounded diameter, or integral-coordinate assumption is imposed.
The workshop's comparison with planar sumsets is motivation rather than an additional hypothesis on $A$.
\subsection{Short English statement}
Determine how small the sumset of $n$ points in three-dimensional space can be when no five points lie in a plane.
\subsection{Sources}
\begin{itemize}
\item[S1] ULAM.AI and the UnsolvedMath Contributors, supplied problems.json, record 20000897 (AIM-COMBINATORICS-0022), AIM Workshop Problem Lists. Catalogue identity and source transcription; CC BY 4.0.
\url{https://huggingface.co/datasets/ulamai/UnsolvedMath}.
\item[S2] American Institute of Mathematics, Additive combinatorics and its applications, item 2.05. Original workshop question underlying AIM-COMBINATORICS-0022.
\url{http://aimpl.org/addcombapp/2/}.
\end{itemize}
\end{document}
